\begin{errata}

\noindent AUTOR, \textbf{\imprimirtitulo} \imprimirsubtitulo. nº de páginas. \imprimirtipotrabalho - 
\imprimirdepartamento, \imprimirinstituicao, \imprimirlocal, \imprimirano.

\begin{table}[!ht]
\centering
\begin{tabular}{|c|c|c|c|} \hline
\textbf{Página} & \textbf{Linha} & \textbf{Onde se lê} & \textbf{Leia-se} \\ \hline
16 & 10 & &  \\ \hline
\end{tabular}
\end{table}


Este elemento pré-textual é opcional e deve, geralmente, ser inserido no texto após o trabalho já ter sido impresso, seguindo a ficha catalográfica ou a folha de rosto. Deve conter a referência do trabalho e o texto da errata com a indicação da página e linha \cite{NBR14724:2011}.

\end{errata}